%% ============================================================
%% Clustering-Guided Contrastive Learning for Clinically Aware
%% Chest X-ray Image--Report Retrieval
%% Compile: xelatex paper.tex  (or pdflatex with T5 + vntex)
%% ============================================================
\documentclass[12pt, a4paper]{article}

% -------- Encoding / Font --------
\usepackage[utf8]{inputenc}
\usepackage[T5]{fontenc}
% Babel: last language listed = main language
\usepackage[english,vietnamese]{babel}

% -------- Maths --------
\usepackage{amsmath, amssymb, amsthm}

% -------- Layout --------
\usepackage[top=2.5cm, bottom=2.5cm, left=3cm, right=2.5cm]{geometry}
\usepackage{setspace}
\onehalfspacing
\usepackage{parskip}
\setlength{\parindent}{1.2em}

% -------- Tables & Figures --------
\usepackage{booktabs, multirow, tabularx, makecell}
\usepackage{graphicx}
\usepackage{caption, subcaption}
\usepackage{float}
\usepackage{array}
\newcolumntype{C}[1]{>{\centering\arraybackslash}m{#1}}
\newcolumntype{L}[1]{>{\raggedright\arraybackslash}m{#1}}

% -------- Algorithm --------
\usepackage{algorithm}
\usepackage{algpseudocode}

% -------- Hyperref --------
\usepackage[colorlinks=true, linkcolor=blue!70!black,
            citecolor=green!50!black, urlcolor=blue]{hyperref}

% -------- References --------
\usepackage{natbib}
\bibliographystyle{abbrvnat}
\setcitestyle{authoryear, open={(}, close={)}}

% -------- Misc --------
\usepackage{xcolor}
\usepackage{enumitem}
\usepackage{microtype}
\usepackage{fancyhdr}
\pagestyle{fancy}
\fancyhf{}
\fancyhead[L]{\small\textit{Học đối sánh có hướng dẫn cụm bệnh cho truy hồi ảnh--báo cáo X-quang ngực}}
\fancyhead[R]{\small\thepage}
\renewcommand{\headrulewidth}{0.4pt}

% -------- Custom commands --------
\newcommand{\Rcal}{\mathcal{R}}
\newcommand{\calL}{\mathcal{L}}
\newcommand{\vect}[1]{\mathbf{#1}}
\newcommand{\norm}[1]{\left\lVert#1\right\rVert}
\newcommand{\ie}{\textit{i.e.}}
\newcommand{\eg}{\textit{e.g.}}
\newcommand{\etal}{\textit{et al.}}
\newcommand{\R}[1]{R@{#1}}

% =========================================================
\begin{document}
% =========================================================

%% ---------------------------------------------------------
%% TITLE
%% ---------------------------------------------------------
\begin{titlepage}
  \centering
  \vspace*{1.5cm}
  {\Large\bfseries
    Học Đối Sánh Có Hướng Dẫn Cụm Bệnh\\
    Cho Truy Hồi Ảnh--Báo Cáo X-quang Ngực\\[0.6em]
    \normalsize\normalfont\itshape
    Clustering-Guided Contrastive Learning for\\
    Clinically Aware Chest X-ray Image--Report Retrieval
  \par}

  \vspace{2cm}
  {\large Diep Trung Nam}\\[0.3em]
  {\normalsize Khoa Công nghệ Thông tin, Trường Đại học \ldots \\[0.2em]
    \texttt{dieptrungnam123@gmail.com}}

  \vspace{2cm}
  {\normalsize \today}

  \vspace{3cm}
  \hrule
  \vspace{0.5em}
  \begin{abstract}
    \noindent\textbf{Tóm tắt (Vietnamese Abstract).}\;
    Truy hồi ảnh--báo cáo X-quang ngực theo hướng học đối sánh đa phương thức
    (\textit{contrastive learning}) thường gặp vấn đề \textit{âm tính sai}
    (\textit{false negative}): các ảnh và báo cáo của bệnh nhân khác nhau nhưng
    cùng nhóm bệnh lý bị xử lý như các cặp không liên quan, gây ra áp lực gradient
    sai lệch và làm suy giảm khả năng truy hồi theo ngữ nghĩa lâm sàng.
    Bài báo này đề xuất khung học đối sánh có hướng dẫn cụm bệnh
    (\textit{clustering-guided contrastive learning framework}) nhằm giảm thiểu
    vấn đề trên. Mô hình kết hợp SwinV2-Base làm bộ mã hóa ảnh và
    Bio\_ClinicalBERT làm bộ mã hóa báo cáo, ánh xạ cả hai về không gian nhúng
    512 chiều dùng chung. Quá trình huấn luyện áp dụng ba giai đoạn: (1) học
    khớp chính xác cặp ảnh--báo cáo bằng InfoNCE đa dương tính, (2) điều chỉnh
    mục tiêu mềm theo độ tương đồng Jaccard có trọng số IDF theo nhóm bệnh lý,
    và (3) kéo các mẫu không bình thường cùng nhóm bệnh lại gần nhau bằng hàm
    mất mát đối sánh có giám sát lâm sàng với lịch curriculum. Trên bộ dữ liệu
    IU-Xray với phân chia cấp độ bệnh nhân cố định, mô hình đạt
    \textit{clinical-valid} mean R@1\,=\,59.86\%,
    \textit{cluster} mean R@1\,=\,69.36\%, và
    \textit{strict} mean R@1\,=\,3.85\% trên tập kiểm tra.
    Trong 377 ca kiểm tra, 64 ca ảnh-sang-văn bản và 93 ca văn bản-sang-ảnh được
    xác định là \textit{clinical rescue}: mô hình trả về kết quả top-1 khác bệnh nhân
    nhưng đúng về nhóm bệnh lý, minh chứng cho hiệu quả của phương pháp đề xuất.

    \medskip
    \noindent\textbf{Từ khoá:} truy hồi ảnh-văn bản y khoa, học đối sánh,
    âm tính sai, cụm bệnh, X-quang ngực, SwinV2, ClinicalBERT.
  \end{abstract}
  \vspace{0.5em}
  \hrule

  \vspace{1.5cm}
  \hrule
  \vspace{0.5em}
  \begin{abstract}
    \noindent\textbf{Abstract (English).}\;
    Cross-modal contrastive learning for chest X-ray image--report retrieval
    commonly suffers from the \emph{false-negative} problem: image--report
    pairs from different patients sharing the same pathological findings are
    treated as hard negatives, introducing incorrect gradient pressure and
    degrading clinically meaningful retrieval.
    We propose a \emph{clustering-guided contrastive learning framework} to
    mitigate this issue. Our model pairs a SwinV2-Base image encoder with a
    Bio\_ClinicalBERT text encoder, projecting both modalities into a shared
    512-dimensional embedding space via deep MLP heads. Training proceeds
    in three phases: (1) strict image--report alignment with multi-positive
    InfoNCE, (2) soft-target blending using IDF-weighted Jaccard pathology
    similarity, and (3) a clinical supervised contrastive loss with a
    curriculum schedule that pulls same-disease, non-normal pairs together
    without compromising exact retrieval.
    On the IU-Xray dataset under a fixed patient-stratified split,
    the proposed model (v8, 150 epochs, NVIDIA A100) achieves a
    \emph{clinical-valid} mean R@1 of 59.86\%,
    \emph{cluster} mean R@1 of 69.36\%, and
    \emph{strict} mean R@1 of 3.85\% on the test set.
    Qualitative analysis reveals 64 image-to-text and 93 text-to-image
    \emph{clinical rescue} cases in the top-1, where the model retrieves
    a clinically related result from a different patient, demonstrating
    effective false-negative mitigation.

    \medskip
    \noindent\textbf{Keywords:} medical image--text retrieval, contrastive
    learning, false negatives, pathology clustering, chest X-ray, SwinV2,
    ClinicalBERT.
  \end{abstract}
  \vspace{0.5em}
  \hrule
\end{titlepage}

\tableofcontents
\newpage

%% ---------------------------------------------------------
%% 1. GIỚI THIỆU
%% ---------------------------------------------------------
\section{Giới Thiệu (Introduction)}
\label{sec:intro}

Truy hồi ảnh--báo cáo y khoa (\textit{medical image--text retrieval}) là bài
toán liên kết ảnh chẩn đoán với báo cáo lâm sàng tương ứng trong không gian
nhúng chung. Ứng dụng thực tế bao gồm hỗ trợ bác sĩ tìm ca tương đồng, tổng
hợp báo cáo tự động và xây dựng hệ thống tư vấn lâm sàng~\citep{jing2018automatic}.

Học đối sánh đa phương thức (\textit{multimodal contrastive learning}) theo phong
cách CLIP~\citep{radford2021learning} đã trở thành nền tảng chủ đạo cho bài toán
này. Các phương pháp như ConVIRT~\citep{zhang2022contrastive} và
GLoRIA~\citep{huang2021gloria} huấn luyện mô hình sao cho các cặp ảnh--báo cáo
\emph{cùng bệnh nhân} (\textit{positive pairs}) được đẩy lại gần nhau, trong khi
các cặp \emph{khác bệnh nhân} (\textit{negative pairs}) bị đẩy xa.

\textbf{Vấn đề âm tính sai (False Negatives).}  Trong y khoa, bệnh nhân khác
nhau thường có cùng biểu hiện bệnh lý. Hình~\ref{fig:false_negative_motivation}
minh họa tình huống này: ảnh X-quang và báo cáo của hai bệnh nhân độc lập nhưng
cùng mắc \textit{Atelectasis} lại bị hàm mất mát đối sánh coi là cặp âm tính và
phạt chúng ra xa nhau. Điều này tạo ra gradient sai lệch và làm hỏng khả năng truy
hồi theo ngữ nghĩa lâm sàng của mô hình.

\begin{figure}[H]
  \centering
  \fbox{\parbox{0.85\textwidth}{
    \centering
    \textbf{Minh họa vấn đề âm tính sai}\\[0.5em]
    \small
    \begin{tabular}{cc}
      Bệnh nhân A: Atelectasis & Bệnh nhân B: Atelectasis \\
      \textit{(Ảnh + Báo cáo A)} & \textit{(Ảnh + Báo cáo B)} \\[0.3em]
      \multicolumn{2}{c}{$\xrightarrow{\text{InfoNCE chuẩn}}$ Coi là negative $\Rightarrow$ Đẩy ra xa} \\[0.3em]
      \multicolumn{2}{c}{$\xrightarrow{\text{Phương pháp đề xuất}}$ Coi là soft positive $\Rightarrow$ Không phạt}
    \end{tabular}
  }}
  \caption{Vấn đề âm tính sai trong học đối sánh y khoa. Hai bệnh nhân khác nhau
    nhưng cùng nhóm bệnh lý bị InfoNCE chuẩn coi là negative. Phương pháp đề xuất
    gán nhãn soft positive dựa trên mức độ tương đồng bệnh lý.}
  \label{fig:false_negative_motivation}
\end{figure}

\textbf{Đóng góp chính của bài báo này:}
\begin{enumerate}[label=(\roman*)]
  \item Xây dựng pipeline truy hồi cấp độ study (\textit{study-level})
    cho bộ dữ liệu IU-Xray, kết hợp tổng hợp đa view bằng
    \textit{attention pooling} (frontal + lateral).
  \item Thiết kế hàm mất mát cụm bệnh theo phân cấp
    (\textit{hierarchical clustering loss}) dùng Jaccard có trọng số IDF
    để tạo nhãn mềm dựa trên mức độ tương đồng bệnh lý.
  \item Đề xuất hàm mất mát đối sánh có giám sát lâm sàng
    (\textit{clinical supervised contrastive loss}) với lịch curriculum,
    kéo các cặp cùng bệnh khác bệnh nhân lại gần nhau mà không làm suy
    yếu truy hồi chính xác.
  \item Xây dựng bộ chỉ số đánh giá ba cấp độ: \textit{strict},
    \textit{cluster}, và \textit{clinical-valid}, trong đó
    \textit{clinical-valid} R@K là chỉ số chính đo hiệu quả giảm thiểu
    âm tính sai.
\end{enumerate}

Phần còn lại của bài báo được tổ chức như sau: Mục~\ref{sec:related} tổng hợp
các công trình liên quan; Mục~\ref{sec:method} mô tả phương pháp đề xuất;
Mục~\ref{sec:exp} trình bày thiết lập thực nghiệm; Mục~\ref{sec:results} báo
cáo kết quả; Mục~\ref{sec:discussion} thảo luận; Mục~\ref{sec:conclusion} kết
luận.

%% ---------------------------------------------------------
%% 2. CÔNG TRÌNH LIÊN QUAN
%% ---------------------------------------------------------
\section{Công Trình Liên Quan (Related Work)}
\label{sec:related}

\subsection{Học Đối Sánh Đa Phương Thức}

CLIP~\citep{radford2021learning} khởi đầu xu hướng huấn luyện ảnh và văn bản
trong không gian nhúng chung bằng hàm InfoNCE~\citep{oord2018representation} trên
400 triệu cặp ảnh--chú thích web. Với dữ liệu y khoa, ConVIRT~\citep{zhang2022contrastive}
là công trình tiên phong áp dụng học đối sánh cho ảnh X-quang và báo cáo lâm sàng.
GLoRIA~\citep{huang2021gloria} mở rộng bằng cách thêm học cục bộ vùng chú ý
(\textit{local attention}). BiomedCLIP~\citep{zhang2023biomedclip} huấn luyện
trên 15 triệu cặp ảnh--văn bản sinh y học quy mô lớn từ PMC.

\subsection{Vấn Đề Âm Tính Sai}

Vấn đề âm tính sai đã được nghiên cứu trong học đối sánh tổng quát.
\citet{khosla2020supervised} đề xuất \textit{Supervised Contrastive Loss}
(SupCon) sử dụng nhãn lớp để xác định nhiều dương tính trong mỗi batch.
\citet{chuang2020debiased} đề xuất ước lượng phân phối âm tính không thiên lệch.
\citet{robinson2021contrastive} nghiên cứu khai thác âm tính khó (\textit{hard
negative mining}) kết hợp tránh âm tính sai.

Trong y khoa, vấn đề này đặc biệt nghiêm trọng do: (1) nhiều bệnh nhân chia sẻ
cùng bệnh lý, (2) dữ liệu có xu hướng mất cân bằng nhãn nặng (\eg{} IU-Xray
có 57.88\% ca bình thường), và (3) ngữ nghĩa bệnh lý phụ thuộc vào cả biểu hiện
hình ảnh lẫn mô tả ngôn ngữ lâm sàng.

\subsection{Bộ Mã Hóa Y Khoa}

SwinV2~\citep{liu2022swin} mở rộng \textit{Swin Transformer} với kỹ thuật
rescaling hậu lớp chuẩn hóa và bias tương đối liên tục (\textit{continuous
relative position bias}), cho phép fine-tune từ độ phân giải thấp lên cao.
Bio\_ClinicalBERT~\citep{alsentzer2019publicly} là mô hình ngôn ngữ BERT
được tiền huấn luyện trên ghi chú bệnh viện MIMIC-III, phù hợp với đặc thù
ngôn ngữ lâm sàng.

\subsection{Truy Hồi Ảnh--Báo Cáo X-quang Ngực}

Bộ dữ liệu IU-Xray~\citep{demner2016preparing} là bộ dữ liệu chuẩn
(\textit{benchmark}) được sử dụng phổ biến nhất cho bài toán tổng hợp báo cáo
và truy hồi X-quang ngực. Nhiều nghiên cứu về truy hồi trên IU-Xray dùng chỉ
số strict R@K với ground-truth là cặp ảnh--báo cáo của cùng bệnh nhân, chưa
xét đến truy hồi theo tương đồng lâm sàng giữa các bệnh nhân khác nhau.
Bài báo này đề xuất chỉ số \textit{clinical-valid} R@K để lấp đầy khoảng
trống này.

%% ---------------------------------------------------------
%% 3. PHƯƠNG PHÁP ĐỀ XUẤT
%% ---------------------------------------------------------
\section{Phương Pháp Đề Xuất (Proposed Method)}
\label{sec:method}

Hình~\ref{fig:architecture} minh họa tổng thể kiến trúc và luồng huấn luyện.

\begin{figure}[H]
  \centering
  \fbox{\parbox{0.95\textwidth}{
    \centering
    \textbf{Kiến trúc tổng thể}\\[1em]
    \small
    \begin{tabular}{ccc}
      \makecell{\textbf{Image Side}\\Frontal (384$\times$384)\\Lateral (384$\times$384)} &
      $\xrightarrow{\text{SwinV2-Base}}$ &
      \makecell{$\vect{v}_f \in \mathbb{R}^{1024}$\\$\vect{v}_l \in \mathbb{R}^{1024}$}
    \end{tabular}
    \vspace{0.5em}\\
    $\downarrow$ \textit{Projection MLP + Attention Pooling}\\
    $\vect{z}_I \in \mathbb{R}^{512}$\\[1em]
    \begin{tabular}{ccc}
      \makecell{\textbf{Text Side}\\Report (max 128 tokens)} &
      $\xrightarrow{\text{Bio\_ClinicalBERT}}$ &
      $\vect{u} \in \mathbb{R}^{768}$
    \end{tabular}
    \vspace{0.5em}\\
    $\downarrow$ \textit{Projection MLP}\\
    $\vect{z}_T \in \mathbb{R}^{512}$\\[1em]
    \textbf{Similarity:} $s_{ij} = e^{\tau} \cdot \vect{z}_I^{(i)\top} \vect{z}_T^{(j)}$
    \quad (Learnable $\tau$, $e^\tau \in [1, 100]$)\\[1em]
    \textbf{Loss:} $\calL = \calL_\text{base} + w_c\,\calL_\text{clinical} + w_a\,\calL_\text{aux}$
  }}
  \caption{Kiến trúc tổng thể của mô hình đề xuất. Hai modality ảnh và văn bản
    được chiếu về không gian nhúng 512 chiều và tính độ tương đồng cosine
    thông qua nhiệt độ học được. Hàm mất mát kết hợp ba thành phần.}
  \label{fig:architecture}
\end{figure}

\subsection{Biểu Diễn Dữ Liệu}
\label{sec:data_rep}

Mỗi mẫu dữ liệu trong IU-Xray tương ứng với một \textit{study} (lần khám) của
một bệnh nhân, gồm:
\begin{itemize}
  \item \textbf{Ảnh:} tối đa một ảnh frontal và một ảnh lateral, kích thước
    $384 \times 384$ pixel sau khi padding để giữ toàn bộ trường nhìn.
  \item \textbf{Văn bản:} báo cáo sau xử lý, ghép từ hai trường \texttt{findings}
    và \texttt{impression}.
  \item \textbf{Nhãn bệnh lý:} vector $\vect{d} \in \{0,1\}^{12}$ biểu diễn
    12 nhóm bệnh lý theo CheXpert~\citep{irvin2019chexpert}, trích xuất bằng
    quy tắc \textit{keyword matching} xét phủ định (\textit{negation-aware}).
    Vector này không bao gồm ``No Finding'' và gộp
    ``Enlarged Cardiomediastinum'' vào ``Cardiomegaly'' theo phép max.
\end{itemize}

\textbf{Định nghĩa bình thường (Normal):}
Sample được coi là \textit{bình thường} (\textit{normal}) khi
$\sum_k d_k = 0$, tức không có nhãn bệnh lý nào bằng 1.

\subsection{Kiến Trúc Mô Hình}
\label{sec:model}

\subsubsection{Bộ Mã Hóa Ảnh (Image Encoder)}

Mô hình sử dụng \textbf{SwinV2-Base}~\citep{liu2022swin} phiên bản
\texttt{microsoft/swinv2-base-patch4-window12to24-192to384-22kto1k-ft}, được
tiền huấn luyện trên ImageNet-22K và fine-tune từ độ phân giải 192 lên 384.
Đầu vào là ảnh $384 \times 384 \times 3$; đầu ra là \texttt{pooler\_output}
có chiều $\mathbb{R}^{1024}$. Gradient checkpointing được bật để tiết kiệm
VRAM.

\subsubsection{Bộ Mã Hóa Văn Bản (Text Encoder)}

Mô hình sử dụng \textbf{Bio\_ClinicalBERT}~\citep{alsentzer2019publicly}
(\texttt{emilyalsentzer/Bio\_ClinicalBERT}), được tiền huấn luyện trên ghi chú
bệnh viện MIMIC-III. Báo cáo được cắt tối đa 128 token; embedding đầu ra là
token \texttt{[CLS]} từ lớp ẩn cuối, chiều $\mathbb{R}^{768}$.

\subsubsection{Đầu Chiếu (Projection Head)}

Cả hai modality đều đi qua MLP 2 lớp với cùng cấu trúc:
\begin{equation}
  \vect{z} = \frac{\text{MLP}(\vect{h})}{\norm{\text{MLP}(\vect{h})}_2},
  \label{eq:proj}
\end{equation}
trong đó MLP được định nghĩa:
\begin{equation}
  \text{MLP}(\vect{h}) = \vect{W}_2 \cdot \text{Dropout}\bigl(
    \text{GELU}(\text{BN}(\vect{W}_1 \vect{h} + \vect{b}_1))
  \bigr) + \vect{b}_2,
\end{equation}
với $\vect{W}_1 \in \mathbb{R}^{1024 \times d_\text{in}}$,
$\vect{W}_2 \in \mathbb{R}^{512 \times 1024}$,
Dropout tỉ lệ 0.2, và BN là Batch Normalization.
Đầu ra $\vect{z} \in \mathbb{R}^{512}$ được chuẩn hóa L2.

Độ tương đồng giữa cặp image--text được tính bằng tích vô hướng có nhiệt độ:
\begin{equation}
  s_{ij} = e^{\tau} \cdot \vect{z}_I^{(i)\top} \vect{z}_T^{(j)},
  \label{eq:sim}
\end{equation}
trong đó $\tau$ là tham số học được (\textit{logit\_scale}), được kẹp sao cho
$e^{\tau} \in [1, 100]$.

\subsubsection{Tổng Hợp Đa View (Multi-View Attention Pooling)}

Mỗi view (frontal/lateral) được mã hóa riêng bởi SwinV2 và projection head.
Sau đó, một \textit{view-type embedding} $\vect{e}_\text{type} \in \mathbb{R}^{512}$
(4 loại: frontal, lateral, other, pad) được cộng vào embedding của view tương ứng.
Embedding study-level được tính bằng \textit{attention pooling}:
\begin{equation}
  \alpha_k = \frac{\exp\bigl(a(\vect{z}_k + \vect{e}_k^\text{type})\bigr)}
                  {\sum_{k'} \exp\bigl(a(\vect{z}_{k'} + \vect{e}_{k'}^\text{type})\bigr)},
  \label{eq:attn_weight}
\end{equation}
\begin{equation}
  \vect{z}_I = \sum_{k=1}^{K} \alpha_k \, \vect{z}_k,
  \label{eq:study_emb}
\end{equation}
trong đó $a(\cdot) = \vect{w}_2^\top \text{GELU}(\vect{W}_1 \cdot)$ là mạng
attention nhỏ với $\vect{W}_1 \in \mathbb{R}^{256 \times 512}$ và
$\vect{w}_2 \in \mathbb{R}^{256}$; $K \leq 2$ là số view hợp lệ
(bị che bởi \textit{view\_mask}). Thiết kế này cho phép mô hình tự học
tầm quan trọng của từng view thay vì dùng hợp trung bình cứng.

\subsubsection{Đầu Phụ Bệnh Lý (Auxiliary Pathology Heads)}

Ngoài đầu chiếu chính, mô hình có bốn đầu phụ (MLP một lớp):
\begin{itemize}
  \item \texttt{img\_pathology\_head}: dự đoán 12 nhãn bệnh lý từ $\vect{z}_I$.
  \item \texttt{img\_normal\_head}: dự đoán nhị phân bình thường/bất thường.
  \item \texttt{txt\_pathology\_head} và \texttt{txt\_normal\_head}: tương tự từ $\vect{z}_T$.
\end{itemize}
Các đầu phụ ép embedding giữ thông tin bệnh lý, nhưng trọng số nhỏ
($w_a = 0.02$) để không lấn át mục tiêu truy hồi.

\subsection{Hàm Mất Mát Theo Phân Cấp Cụm Bệnh}
\label{sec:loss}

\subsubsection{Giai Đoạn 1: MultiPositive InfoNCE (Epoch 1--11)}

Trước khi bật clustering, hàm mất mát chính là InfoNCE đa dương tính:
\begin{equation}
  \calL_\text{InfoNCE} = -\frac{1}{N}\sum_{i=1}^{N}
    \log \frac{\exp(s_{ii})}
              {\sum_{j=1}^{N} \exp(s_{ij})},
  \label{eq:infonce}
\end{equation}
trong đó $s_{ij}$ theo công thức~\eqref{eq:sim} đã tích hợp nhiệt độ học được,
nên không cần chia thêm. Hàm mất mát tính theo cả hai chiều i2t và t2i;
ground-truth dương tính là cặp cùng \texttt{patient\_id}.

\subsubsection{Giai Đoạn 2: Hierarchical Clustering Loss (Epoch 12+)}
\label{sec:cluster_loss}

Từ epoch thứ 12, ground-truth được điều chỉnh thành nhãn mềm
(\textit{soft target}) kết hợp giữa khớp chính xác và tương đồng bệnh lý:
\begin{equation}
  t_{ij} = (1 - \alpha) \cdot \mathbb{1}[p_i = p_j]
           + \alpha \cdot S^\text{IDF}_{ij},
  \label{eq:soft_target}
\end{equation}
với $\alpha = 0.35$ và $S^\text{IDF}_{ij}$ là độ tương đồng Jaccard có trọng
số IDF giữa hai vector bệnh lý $\vect{d}_i$ và $\vect{d}_j$:

\begin{equation}
  \boxed{
  S^\text{IDF}(A, B) =
    \frac{\displaystyle\sum_{k=1}^{12} w_k \cdot \min(A_k, B_k)}
         {\displaystyle\sum_{k=1}^{12} w_k \cdot \max(A_k, B_k)},
  }
  \label{eq:jaccard_idf}
\end{equation}
trong đó trọng số IDF được tính:
\begin{equation}
  w_k = \frac{\log(1/f_k)}{\sum_{k'} \log(1/f_{k'})},
  \label{eq:idf_weight}
\end{equation}
với $f_k$ là tần suất xuất hiện của bệnh $k$ trong tập huấn luyện.
Thiết kế này ưu tiên bệnh hiếm gặp hơn: \eg{} Pneumothorax ($f = 2.03\%$) và
Fracture ($f = 1.82\%$) có $w_k$ lớn hơn Cardiomegaly ($f = 9.00\%$).

\textbf{Xử lý nhóm bình thường:}
Khi $\sum_k d_k^{(i)} = 0$ và $\sum_k d_k^{(j)} = 0$ (cả hai đều bình thường),
Jaccard không xác định được. Vì vậy, công thức~\eqref{eq:soft_target} thêm
phần thưởng cố định:
\begin{equation}
  S^\text{IDF}_{ij} \leftarrow S^\text{IDF}_{ij} + 0.4 \cdot
    \mathbb{1}[\text{normal}(i)] \cdot \mathbb{1}[\text{normal}(j)],
  \label{eq:normal_bonus}
\end{equation}
giúp nhóm các ca bình thường lại gần nhau trong không gian nhúng.

\subsubsection{Hàm Mất Mát Đối Sánh Có Giám Sát Lâm Sàng (Epoch 15+)}
\label{sec:clinical_loss}

Đây là đóng góp cốt lõi nhắm trực tiếp vào vấn đề âm tính sai:

\textbf{Mặt nạ lâm sàng (Clinical mask):}
\begin{equation}
  M^\text{clin}_{ij} =
    \underbrace{\mathbb{1}[\vect{d}_i \cdot \vect{d}_j > 0]}_{\text{có disease overlap}}
    \cdot
    \underbrace{\mathbb{1}[i \neq j]}_{\text{khác sample}}
    \cdot
    \underbrace{\mathbb{1}[\text{non-normal}(i)]}_{}
    \cdot
    \underbrace{\mathbb{1}[\text{non-normal}(j)]}_{}
  \label{eq:clinical_mask}
\end{equation}
Mặt nạ này xác định các cặp bệnh nhân khác nhau, cùng có ít nhất một
bệnh lý chung và đều không bình thường.

\textbf{Hàm mất mát:}
\begin{equation}
  \calL_\text{clin} = -\frac{1}{|\mathcal{Q}|} \sum_{i \in \mathcal{Q}}
    \frac{1}{|P_i|} \sum_{j \in P_i}
    \log \frac{\exp(s_{ij})}{\sum_{k=1}^{N} \exp(s_{ik})},
  \label{eq:clinical_loss}
\end{equation}
trong đó $\mathcal{Q} = \{i : |P_i| > 0\}$ là tập query bất thường có ít nhất
một positive lâm sàng, và $P_i = \{j : M^\text{clin}_{ij} = 1\}$. Hàm mất mát
được tính theo cả hai chiều image$\to$text và text$\to$image.

\subsubsection{Hàm Mất Mát Phụ (Auxiliary Loss)}

\begin{equation}
  \calL_\text{aux} = \frac{1}{4}\bigl[
    \text{BCE}(\hat{\vect{d}}_I, \vect{d}) + \text{BCE}(\hat{n}_I, n)
    + \text{BCE}(\hat{\vect{d}}_T, \vect{d}) + \text{BCE}(\hat{n}_T, n)
  \bigr]
  + 0.25 \cdot \bigl[\text{MSE}(\sigma(\hat{\vect{d}}_I), \sigma(\hat{\vect{d}}_T))
    + \text{MSE}(\sigma(\hat{n}_I), \sigma(\hat{n}_T))\bigr],
  \label{eq:aux_loss}
\end{equation}
trong đó $\hat{\vect{d}}$ là dự đoán bệnh lý, $\hat{n}$ là dự đoán bình thường,
và $\sigma(\cdot)$ là sigmoid. Thành phần MSE ép hai modality nhất quán về dự
đoán bệnh lý.

\textbf{Hàm mất mát tổng:}
\begin{equation}
  \calL = \calL_\text{base} + w_c(e) \cdot \calL_\text{clin} + w_a \cdot \calL_\text{aux},
  \label{eq:total_loss}
\end{equation}
với $w_a = 0.02$ và $w_c(e)$ là lịch curriculum phụ thuộc epoch $e$.

\subsection{Lịch Curriculum Lâm Sàng}
\label{sec:curriculum}

Thay vì bật $\calL_\text{clin}$ mạnh ngay từ đầu---có thể kéo mô hình về cụm
quá rộng trước khi học được khớp chính xác---bài báo này áp dụng lịch học tăng
dần rồi giảm:

\begin{equation}
  w_c(e) =
  \begin{cases}
    0 & e < 15 \\[2pt]
    0.03 + 0.09 \cdot \dfrac{e - 15}{25 - 15} & 15 \leq e \leq 25 \\[6pt]
    0.12 & 26 \leq e \leq 60 \\[2pt]
    0.08 & 61 \leq e \leq 105 \\[2pt]
    0.04 & 106 \leq e \leq 150
  \end{cases}
  \label{eq:curriculum}
\end{equation}

Bảng~\ref{tab:curriculum} tóm tắt ý nghĩa từng giai đoạn.

\begin{table}[H]
  \centering
  \caption{Lịch curriculum lâm sàng và mục tiêu từng giai đoạn.}
  \label{tab:curriculum}
  \begin{tabular}{clcc}
    \toprule
    \textbf{Epoch} & \textbf{Mục tiêu chính} & \textbf{$w_c$} & \textbf{Pha huấn luyện} \\
    \midrule
    1--5   & Chỉ huấn luyện đầu chiếu (backbone đóng băng)   & 0.00 & Warm-up \\
    6--11  & MultiPositive InfoNCE, mở toàn bộ backbone       & 0.00 & Strict alignment \\
    12--14 & Bật clustering soft target                       & 0.00 & Clustering phase \\
    15--25 & Tăng dần clinical loss                           & 0.03 $\to$ 0.12 & Clinical ramp \\
    26--60 & Clinical supervision đỉnh                        & 0.12 & Clinical peak \\
    61--105& Giảm dần clinical để hồi phục strict             & 0.08 & Clinical decay \\
    106--150& Giai đoạn cuối, giảm thêm clinical              & 0.04 & Fine-tune strict \\
    \bottomrule
  \end{tabular}
\end{table}

Kết quả thực nghiệm (Mục~\ref{sec:results}) xác nhận tính hợp lý của lịch này:
strict R@1 đạt đỉnh tại epoch 120 sau khi clinical weight giảm, trong khi
clinical-valid đạt đỉnh tại epoch 90 khi clinical weight vẫn còn cao.

%% ---------------------------------------------------------
%% 4. THỰC NGHIỆM
%% ---------------------------------------------------------
\section{Thực Nghiệm (Experiments)}
\label{sec:exp}

\subsection{Bộ Dữ Liệu}
\label{sec:dataset}

Bộ dữ liệu \textbf{IU-Xray}~\citep{demner2016preparing} (Indiana University
Chest X-ray Collection) gồm ảnh X-quang ngực chụp từ hệ thống PACS bệnh viện
Indiana, kèm theo báo cáo lâm sàng tiếng Anh đã ẩn danh.

\textbf{Tiền xử lý:}
\begin{itemize}
  \item Ảnh được chuyển sang RGB, padding thành hình vuông bằng nền đen, rồi
    lưu ở kích thước $384 \times 384$ để tương thích với SwinV2-Base.
    Padding được ưu tiên hơn crop để bảo toàn toàn bộ trường nhìn lâm sàng.
  \item Văn bản: ghép \texttt{findings} + \texttt{impression}, chuẩn hóa
    chữ thường, loại bỏ placeholder và thông tin nhân khẩu, lọc báo cáo
    dưới 10 từ.
  \item Nhãn bệnh lý: 14 nhãn CheXpert được tạo bằng quy tắc keyword matching
    xét phủ định, sau đó rút gọn thành vector 12 chiều.
\end{itemize}

Bảng~\ref{tab:dataset} thống kê dữ liệu sau xử lý.

\begin{table}[H]
  \centering
  \caption{Thống kê bộ dữ liệu IU-Xray sau tiền xử lý.}
  \label{tab:dataset}
  \begin{tabular}{lrr}
    \toprule
    \textbf{Thông số} & \multicolumn{2}{c}{\textbf{Giá trị}} \\
    \midrule
    Tổng số ảnh (PNG, 384$\times$384)    & \multicolumn{2}{r}{7,322} \\
    Số bệnh nhân / study ID duy nhất     & \multicolumn{2}{r}{3,772} \\
    Số báo cáo duy nhất                  & \multicolumn{2}{r}{3,018} \\
    Số từ (min / trung bình / max)       & \multicolumn{2}{r}{10 / 37.27 / 223} \\
    \midrule
    Projection & Frontal & 3,738 \\
               & Lateral & 3,584 \\
    \midrule
    \textbf{Phân chia dữ liệu (seed=42, patient-stratified)} & & \\
    Huấn luyện (train) -- study-level   & \multicolumn{2}{r}{3,018} \\
    Kiểm định (val) -- study-level      & \multicolumn{2}{r}{377}   \\
    Kiểm tra (test) -- study-level      & \multicolumn{2}{r}{377}   \\
    \bottomrule
  \end{tabular}
\end{table}

\textbf{Phân bố nhãn bệnh lý} được trình bày trong Bảng~\ref{tab:label_dist}.

\begin{table}[H]
  \centering
  \caption{Phân bố 14 nhãn CheXpert trong tập dữ liệu IU-Xray (7,322 ảnh).}
  \label{tab:label_dist}
  \begin{tabular}{lrr}
    \toprule
    \textbf{Nhãn bệnh lý} & \textbf{Số lượng} & \textbf{Tỉ lệ (\%)} \\
    \midrule
    No Finding              & 4,238 & 57.88 \\
    Cardiomegaly            &   659 &  9.00 \\
    Lung Lesion             & 1,014 & 13.85 \\
    Lung Opacity            &   942 & 12.87 \\
    Atelectasis             &   732 & 10.00 \\
    Consolidation           &   402 &  5.49 \\
    Pleural Effusion        &   415 &  5.67 \\
    Support Devices         &   224 &  3.06 \\
    Edema                   &   185 &  2.53 \\
    Pneumonia               &   164 &  2.24 \\
    Pneumothorax            &   149 &  2.03 \\
    Fracture                &   133 &  1.82 \\
    Pleural Other           &    65 &  0.89 \\
    Enlarged Cardiomediastinum &  17 &  0.23 \\
    \bottomrule
  \end{tabular}
\end{table}

\subsection{Cài Đặt Thực Nghiệm}
\label{sec:setup}

\textbf{Phần cứng:} NVIDIA A100 40 GB, Python 3.x, PyTorch, HuggingFace
Transformers. Mixed precision (\texttt{torch.amp.autocast + GradScaler}).

Các siêu tham số chính được tóm tắt trong Bảng~\ref{tab:hparams}.

\begin{table}[H]
  \centering
  \caption{Siêu tham số huấn luyện của mô hình v8 (run chính).}
  \label{tab:hparams}
  \begin{tabular}{lc}
    \toprule
    \textbf{Siêu tham số} & \textbf{Giá trị} \\
    \midrule
    Số epoch                     & 150 \\
    Batch size                   & 32 \\
    Gradient accumulation        & 4 \\
    Effective batch size         & 128 \\
    Epoch đóng băng backbone     & 5 \\
    \midrule
    Learning rate (image encoder)  & $1 \times 10^{-5}$ \\
    Learning rate (text encoder)   & $1 \times 10^{-4}$ \\
    Learning rate (projection head)& $2 \times 10^{-4}$ \\
    Weight decay                   & 0.01 \\
    Max gradient norm              & 1.0 \\
    Optimizer                      & AdamW \\
    \midrule
    LR Scheduler (sau unfreeze) & Linear warmup (5\%) + CosineAnnealing \\
    $\eta_\text{min}$           & $10^{-8}$ \\
    \midrule
    $\alpha$ (soft target blend)  & 0.35 \\
    Temperature (cluster)         & 1.0 \\
    $w_a$ (auxiliary weight)      & 0.02 \\
    Seed                          & 42 \\
    \bottomrule
  \end{tabular}
\end{table}

\textbf{Chiến lược checkpoint:}
Mô hình lưu ba loại checkpoint trong quá trình huấn luyện:
(1) \texttt{best.pt}: epoch có \textit{strict} R@1 validation cao nhất;
(2) \texttt{best\_balanced.pt}: epoch tối ưu hóa $\text{strict} + 0.1 \cdot
\text{clinical-valid}$ khi strict $\geq 4.0$\%;
(3) \texttt{best\_clinical\_valid.pt}: epoch có \textit{clinical-valid} R@1
validation cao nhất khi strict $\geq 3.5$\%.
Checkpoint chính được báo cáo là \texttt{best.pt} (epoch 120).

\subsection{Chỉ Số Đánh Giá}
\label{sec:metrics}

Đánh giá ở cấp độ study: mỗi bệnh nhân là một query, embedding study được
tính bằng attention pooling giống lúc huấn luyện.

\textbf{(M1) Strict Retrieval:}
Ground truth $G^\text{strict}_{ij} = \mathbb{1}[p_i = p_j]$, tức chỉ đúng
bệnh nhân mới được tính là hit. Đây là chỉ số khó nhất, đo khớp chính xác.
Báo cáo R@1, R@5, R@10, MRR theo cả hai chiều i2t và t2i.

\textbf{(M2) Cluster Retrieval:}
$G^\text{clus}_{ij} = \mathbb{1}[\vect{d}_i \cdot \vect{d}_j > 0]
+ \mathbb{1}[\text{normal}(i) \wedge \text{normal}(j)]$.
Đo retrieval theo cụm bệnh rộng, bao gồm cả cặp bình thường--bình thường.

\textbf{(M3) Clinical-Valid Retrieval (Chỉ số chính):}
$G^\text{clin}_{ij} = \mathbb{1}[\vect{d}_i \cdot \vect{d}_j > 0]
\wedge \mathbb{1}[\text{non-normal}(i)] \wedge \mathbb{1}[\text{non-normal}(j)]$.
Chỉ tính metric trên các query \textit{non-normal} có ít nhất một positive lâm
sàng trong tập kiểm tra ($n = 147$ query mỗi chiều).
\textit{Đây là chỉ số đo trực tiếp mục tiêu giảm âm tính sai.}

\textbf{(M4) Clinical-All Retrieval:} Tương tự M3 nhưng tính trên tất cả 377
query (kể cả query không có positive lâm sàng). Dùng làm chỉ số phụ.

\textbf{Phân tích định tính (Bucket Analysis):}
Mỗi kết quả top-1 được phân vào một trong bốn nhóm:
\begin{itemize}
  \item \textit{Strict hit}: top-1 đúng bệnh nhân.
  \item \textit{Clinical rescue}: top-1 khác bệnh nhân nhưng đúng clinical mask.
  \item \textit{Cluster rescue}: top-1 đúng cluster mask nhưng không đủ chuẩn clinical.
  \item \textit{Miss}: top-1 không thuộc nhóm nào ở trên.
\end{itemize}

%% ---------------------------------------------------------
%% 5. KẾT QUẢ
%% ---------------------------------------------------------
\section{Kết Quả (Results)}
\label{sec:results}

\subsection{Kết Quả Định Lượng Trên Tập Kiểm Tra}

Bảng~\ref{tab:main_results} trình bày kết quả đầy đủ của
\textbf{v8 (epoch 120, best.pt)} trên tập kiểm tra 377 study.

\begin{table}[H]
  \centering
  \caption{Kết quả truy hồi trên tập kiểm tra IU-Xray.
    \textbf{i2t}: image-to-text; \textbf{t2i}: text-to-image; \textbf{Mean}: trung bình hai chiều.
    Metric chính: \textit{clinical-valid} R@1 (in đậm).
    Checkpoint: \texttt{best.pt}, epoch 120.}
  \label{tab:main_results}
  \renewcommand{\arraystretch}{1.15}
  \begin{tabular}{l ccc ccc ccc}
    \toprule
    & \multicolumn{3}{c}{\textbf{Strict}} &
      \multicolumn{3}{c}{\textbf{Cluster}} &
      \multicolumn{3}{c}{\textbf{Clinical-Valid}} \\
    \cmidrule(lr){2-4} \cmidrule(lr){5-7} \cmidrule(lr){8-10}
    & R@1 & R@5 & R@10 & R@1 & R@5 & R@10 & R@1 & R@5 & R@10 \\
    \midrule
    i2t
      & 3.45 & 10.08 & 17.77
      & 68.17 & 79.58 & 85.15
      & 49.66 & 60.54 & 69.39 \\
    t2i
      & 4.24 & 12.20 & 18.83
      & 70.56 & 97.08 & 99.20
      & 70.07 & 95.24 & 97.96 \\
    \midrule
    \textbf{Mean}
      & 3.85 & 11.14 & 18.30
      & 69.36 & 88.33 & 92.18
      & \textbf{59.86} & \textbf{77.89} & \textbf{83.67} \\
    \bottomrule
  \end{tabular}
  \vspace{0.4em}\\
  \small \textit{MRR (strict) = 9.01\%; clinical-all mean R@1 = 23.34\%;
  clinical-valid: 147/377 query mỗi chiều có positive.}
\end{table}

\textbf{Kết quả tập kiểm định (Validation, epoch 120):}
\begin{itemize}
  \item Strict R@1 (mean) tốt nhất: \textbf{4.91\%} tại epoch 120.
  \item Clinical-valid R@1 (mean) tốt nhất: \textbf{52.11\%} tại epoch 90.
  \item Cluster R@1 (mean) tốt nhất: \textbf{65.25\%} tại epoch 40.
\end{itemize}

\subsection{Phân Tích Định Tính: Bucket Analysis}

Bảng~\ref{tab:bucket} phân tích top-1 của 377 study kiểm tra.

\begin{table}[H]
  \centering
  \caption{Phân tích định tính top-1 trên 377 study kiểm tra. \textit{Clinical
    rescue} là các ca mô hình trả về kết quả khác bệnh nhân nhưng đúng về
    nhóm bệnh lý --- minh chứng trực tiếp cho hiệu quả giảm âm tính sai.}
  \label{tab:bucket}
  \begin{tabular}{lcccc}
    \toprule
    \textbf{Chiều} &
    \textbf{Strict Hit} &
    \textbf{Clinical Rescue} &
    \textbf{Cluster Rescue} &
    \textbf{Miss} \\
    \midrule
    Image $\to$ Text & 13 (3.4\%)  & 64 (17.0\%) & 180 (47.7\%) & 120 (31.8\%) \\
    Text $\to$ Image & 16 (4.2\%)  & 93 (24.7\%) & 157 (41.6\%) & 111 (29.4\%) \\
    \midrule
    \textbf{Tổng}    & \textbf{29} & \textbf{157} & \textbf{337} & \textbf{231} \\
    \bottomrule
  \end{tabular}
\end{table}

Tổng cộng \textbf{157 ca clinical rescue} (64 i2t + 93 t2i) xác nhận mô hình
học được tương đồng lâm sàng vượt qua ranh giới bệnh nhân. Trong 29 ca strict
hit, top-1 đúng hoàn toàn bệnh nhân. 337 ca cluster rescue cho thấy mô hình
nhóm tốt các ca có cùng cụm bệnh rộng. Chỉ 231/754 ca (30.6\%) bị miss hoàn
toàn.

\subsection{So Sánh Với Mô Hình Nền}
\label{sec:comparison}

Bảng~\ref{tab:comparison} so sánh v8 (mô hình đề xuất) với v7 trên cùng
phân chia dữ liệu (seed=42, patient-stratified).

\begin{table}[H]
  \centering
  \caption{So sánh v8 (đề xuất) với v7 (\textit{baseline}) trên tập kiểm tra.
    Cả hai được đánh giá trên cùng phân chia dữ liệu cố định.}
  \label{tab:comparison}
  \begin{tabular}{l cccc}
    \toprule
    \textbf{Mô hình} &
    \textbf{Strict Mean R@1} &
    \textbf{Cluster Mean R@1} &
    \textbf{Clinical-Valid Mean R@1} &
    \textbf{MRR} \\
    \midrule
    V7 (best-balanced checkpoint)
      & 3.32\% & 67.37\% & 59.86\% & -- \\
    \textbf{V8 (best.pt, ep.~120)}
      & \textbf{3.85\%} & \textbf{69.36\%} & \textbf{59.86\%} & \textbf{9.01\%} \\
    \midrule
    \textit{Cải thiện ($\Delta$)}
      & +0.53 pp & +1.99 pp & 0.00 pp & -- \\
    \bottomrule
  \end{tabular}
\end{table}

V8 cải thiện strict mean R@1 thêm \textbf{+0.53 pp} và cluster mean R@1 thêm
\textbf{+1.99 pp} so với v7, trong khi duy trì clinical-valid mean R@1 ở mức
59.86\%. Kết quả này cho thấy lịch curriculum của v8 cân bằng tốt hơn giữa
truy hồi chính xác và truy hồi lâm sàng so với cài đặt của v7.

\subsection{Động Lực Huấn Luyện}

Hình~\ref{fig:training_dynamics} minh họa diễn biến các chỉ số qua 150 epoch.

\begin{figure}[H]
  \centering
  \fbox{\parbox{0.92\textwidth}{
    \centering
    \small
    \begin{tabular}{r l}
      \textbf{Epoch 1--5:}  & Chỉ huấn luyện đầu chiếu (backbone đóng băng). \\
      \textbf{Epoch 6:}     & Mở backbone; strict R@1 bắt đầu tăng mạnh. \\
      \textbf{Epoch 12:}    & Bật clustering loss; cluster R@1 tăng rõ rệt. \\
      \textbf{Epoch 15:}    & Bắt đầu tăng $w_c$; clinical-valid tăng dần. \\
      \textbf{Epoch 40:}    & Cluster R@1 đạt đỉnh validation (65.25\%). \\
      \textbf{Epoch 60:}    & $w_c$ bắt đầu giảm; strict R@1 phục hồi. \\
      \textbf{Epoch 90:}    & Clinical-valid R@1 đạt đỉnh validation (52.11\%). \\
      \textbf{Epoch 120:}   & \textbf{Strict R@1 đạt đỉnh validation (4.91\%)} --- checkpoint best.pt. \\
      \textbf{Epoch 150:}   & Epoch cuối; strict giảm nhẹ so với epoch 120.
    \end{tabular}
  }}
  \caption{Diễn biến các chỉ số trên tập kiểm định qua 150 epoch.
    (Hình ảnh thực tế từ \texttt{history.csv} nên được vẽ bằng matplotlib.)
    Điểm đặc trưng: strict đạt đỉnh muộn (epoch 120) sau khi clinical decay,
    xác nhận hiệu quả của lịch curriculum.}
  \label{fig:training_dynamics}
\end{figure}

Điểm quan trọng cần nhấn mạnh: strict R@1 \textbf{phục hồi rõ ràng} sau khi
$w_c$ bắt đầu giảm tại epoch 60, và đạt đỉnh tại epoch 120. Điều này xác nhận
giả thuyết thiết kế: clinical decay giúp mô hình ``tinh chỉnh lại'' khả năng
khớp chính xác sau giai đoạn học cụm bệnh.

%% ---------------------------------------------------------
%% 6. THẢO LUẬN
%% ---------------------------------------------------------
\section{Thảo Luận (Discussion)}
\label{sec:discussion}

\subsection{Ý Nghĩa Của Clinical-Valid R@1 Cao}

Clinical-valid mean R@1 = 59.86\% trên 147 query (mỗi chiều) có ít nhất một
positive lâm sàng cho thấy mô hình học được tương đồng bệnh lý vượt qua ranh
giới bệnh nhân. Trong thực tế lâm sàng, khả năng truy hồi ca tương đồng bệnh
lý (dù khác bệnh nhân) hữu ích hơn khớp chính xác---bác sĩ cần ca tham chiếu
có biểu hiện giống, không nhất thiết cùng bệnh nhân.

\subsection{Hạn Chế Của Strict R@1}

Strict R@1 = 3.85\% thấp có thể giải thích bởi: (1) IU-Xray nhỏ (3,018 study
huấn luyện) tương đối so với độ phức tạp của bài toán;
(2) nhiều báo cáo IU-Xray có nội dung rất giống nhau, gây khó khăn cho việc
phân biệt chính xác; (3) mô hình đang học đồng thời hai mục tiêu (strict và
clinical), tạo đánh đổi (\textit{trade-off}).

\subsection{Vai Trò Của Jaccard IDF và Healthy Cluster}

Trọng số IDF đảm bảo rằng các bệnh hiếm gặp như Pneumothorax (2.03\%) và
Fracture (1.82\%) đóng góp nhiều hơn vào độ tương đồng so với bệnh phổ biến
như Cardiomegaly (9.00\%). Điều này phù hợp với thực tế lâm sàng: hai ảnh
cùng có Pneumothorax hiếm gặp đáng được xem là tương đồng hơn hai ảnh cùng
có Cardiomegaly phổ biến. Nhóm bình thường (57.88\% dữ liệu) được xử lý
riêng qua phần thưởng 0.4, tránh trường hợp vector zero làm Jaccard không
xác định được.

\subsection{Curriculum Vs. Fixed Clinical Weight}

Lịch curriculum (tăng rồi giảm $w_c$) được thiết kế để giải quyết bài toán
cân bằng: nếu $w_c$ cố định cao, mô hình bị kéo quá mạnh về cụm rộng và
không học được khớp chính xác; nếu $w_c$ cố định thấp, lợi ích từ clinical
supervision không đủ. Kết quả epoch 120 (best strict) sau giai đoạn
decay xác nhận hiệu quả của thiết kế này.

\subsection{Các Thành Phần Không Bật Trong V8}

Để đảm bảo tính minh bạch khoa học, cần ghi rõ: các thành phần
\textit{DualPrototypeBank}, \textit{IntraClusterRankingLoss}, và
\textit{Hard Negative Mining} có trong codebase nhưng bị vô hiệu hóa trong
run v8 chính (\texttt{proto\_start=999}, \texttt{rank\_start=999},
\texttt{mine\_start=999}). Kết quả báo cáo \textbf{không đến từ} các thành
phần này. Chúng là hướng mở rộng cho công việc tương lai.

\subsection{Hạn Chế}

\begin{enumerate}
  \item \textbf{Quy mô dữ liệu:} IU-Xray chỉ có $\sim$3,000 study huấn
    luyện---nhỏ so với các hệ thống CLIP y khoa hiện đại dùng hàng triệu cặp.
  \item \textbf{Nhãn bệnh lý yếu:} Nhãn được tạo bằng keyword matching, không
    phải annotation của bác sĩ. Một số nhãn có thể sai do xử lý phủ định đơn
    giản.
  \item \textbf{Không có baseline từ paper khác:} Chưa so sánh với ConVIRT,
    GLoRIA, BiomedCLIP trên cùng phân chia dữ liệu, do không có cài đặt mã
    nguồn mở phù hợp.
  \item \textbf{Chỉ đánh giá trên IU-Xray:} Khả năng tổng quát hóa sang
    MIMIC-CXR hay các bộ dữ liệu X-quang khác chưa được kiểm chứng.
  \item \textbf{Cluster metric bị ảnh hưởng bởi nhóm bình thường:} Cluster
    R@1 cao một phần do nhóm normal lớn (57.88\%) tạo nhiều true positive
    dễ.
\end{enumerate}

%% ---------------------------------------------------------
%% 7. KẾT LUẬN
%% ---------------------------------------------------------
\section{Kết Luận (Conclusion)}
\label{sec:conclusion}

Bài báo đề xuất \textbf{khung học đối sánh có hướng dẫn cụm bệnh}
(\textit{clustering-guided contrastive learning framework}) cho bài toán truy
hồi ảnh--báo cáo X-quang ngực, với mục tiêu giảm thiểu vấn đề âm tính sai
trong học đối sánh y khoa. Mô hình kết hợp SwinV2-Base và Bio\_ClinicalBERT,
tổng hợp đa view bằng attention pooling, và huấn luyện theo ba giai đoạn với
lịch curriculum lâm sàng.

Trên IU-Xray với phân chia patient-stratified cố định (seed=42):
\begin{itemize}
  \item \textit{Clinical-valid} mean R@1 = \textbf{59.86\%}
    (chỉ số chính, 147 query/chiều).
  \item \textit{Cluster} mean R@1 = \textbf{69.36\%}.
  \item \textit{Strict} mean R@1 = \textbf{3.85\%}; MRR = 9.01\%.
  \item \textbf{157 ca clinical rescue} trong top-1 (64 i2t + 93 t2i)
    từ 377 study kiểm tra.
\end{itemize}

Kết quả cho thấy phương pháp đề xuất học được tương đồng lâm sàng hiệu quả,
phù hợp với mục tiêu giảm âm tính sai. Strict retrieval còn là hạn chế cần
tiếp tục cải thiện, đặc biệt bằng cách thêm dữ liệu lớn hơn hoặc ablation
có hệ thống hơn. Hướng nghiên cứu tương lai bao gồm: (1) ablation loại từng
thành phần (IDF Jaccard, healthy cluster, curriculum schedule), (2) tích hợp
hard negative mining, (3) đánh giá trên MIMIC-CXR quy mô lớn hơn, và
(4) so sánh với baseline đã published trên cùng phân chia dữ liệu.

%% ---------------------------------------------------------
%% ACKNOWLEDGEMENTS
%% ---------------------------------------------------------
\section*{Lời Cảm Ơn}

Tác giả cảm ơn giáo viên hướng dẫn và các thành viên nhóm nghiên cứu đã có
những góp ý quý báu. Thực nghiệm được thực hiện trên NVIDIA A100 thông qua
Google Colaboratory và nền tảng Kaggle.

%% ---------------------------------------------------------
%% REFERENCES
%% ---------------------------------------------------------
\begin{thebibliography}{99}

\bibitem[Alsentzer et~al.(2019)]{alsentzer2019publicly}
Alsentzer, E., Murphy, J.~R., Boag, W., Weng, W.-H., Jin, D., Naumann, T., and
  McDermott, M. (2019).
\newblock Publicly available clinical {BERT} embeddings.
\newblock In \textit{Proceedings of the 2nd Clinical Natural Language
  Processing Workshop}, pp.~72--78. Association for Computational Linguistics.

\bibitem[Chuang et~al.(2020)]{chuang2020debiased}
Chuang, C.-Y., Robinson, J., Lin, Y.-C., Torralba, A., and Jegelka, S. (2020).
\newblock Debiased contrastive learning.
\newblock In \textit{Advances in Neural Information Processing Systems
  (NeurIPS)}, Vol.~33, pp.~8765--8775.

\bibitem[Demner-Fushman et~al.(2016)]{demner2016preparing}
Demner-Fushman, D., Kohli, M.~D., Rosenman, M.~B., Shooshan, S.~E., Rodriguez,
  L., Antani, S., Thoma, G.~R., and McDonald, C.~J. (2016).
\newblock Preparing a collection of radiology examinations for distribution and
  retrieval.
\newblock \textit{Journal of the American Medical Informatics Association
  (JAMIA)}, 23(2):304--310.

\bibitem[Huang et~al.(2021)]{huang2021gloria}
Huang, S.-C., Shen, L., Lungren, M.~P., and Yeung, S. (2021).
\newblock {GloRIA}: A multimodal global-local representation learning framework
  for label-efficient medical image recognition.
\newblock In \textit{Proceedings of the IEEE/CVF International Conference on
  Computer Vision (ICCV)}, pp.~3942--3951.

\bibitem[Irvin et~al.(2019)]{irvin2019chexpert}
Irvin, J., Rajpurkar, P., Ko, M., Yu, Y., Ciurea-Ilcus, S., Chute, C.,
  Marklund, H., Haghgoo, B., Ball, R., Shpanskaya, K., et~al. (2019).
\newblock {CheXpert}: A large chest radiograph dataset with uncertainty labels
  and expert comparison.
\newblock In \textit{Proceedings of the AAAI Conference on Artificial
  Intelligence}, Vol.~33, pp.~590--597.

\bibitem[Jing et~al.(2018)]{jing2018automatic}
Jing, B., Xie, P., and Xing, E. (2018).
\newblock On the automatic generation of medical imaging reports.
\newblock In \textit{Proceedings of the 56th Annual Meeting of the Association
  for Computational Linguistics (ACL)}, pp.~2577--2586.

\bibitem[Khosla et~al.(2020)]{khosla2020supervised}
Khosla, P., Tian, Y., Wang, X., Liu, C., Isola, P., Krishnamurthy, A., and
  Tian, Y. (2020).
\newblock Supervised contrastive learning.
\newblock In \textit{Advances in Neural Information Processing Systems
  (NeurIPS)}, Vol.~33, pp.~18661--18673.

\bibitem[Liu et~al.(2022)]{liu2022swin}
Liu, Z., Hu, H., Lin, Y., Yao, Z., Xie, Z., Wei, Y., Ning, J., Cao, Y., Zhang,
  Z., Dong, L., et~al. (2022).
\newblock Swin transformer {V2}: Scaling up capacity and resolution.
\newblock In \textit{Proceedings of the IEEE/CVF Conference on Computer Vision
  and Pattern Recognition (CVPR)}, pp.~12009--12019.

\bibitem[Oord et~al.(2018)]{oord2018representation}
Oord, A. v.~d., Li, Y., and Vinyals, O. (2018).
\newblock Representation learning with contrastive predictive coding.
\newblock \textit{arXiv preprint arXiv:1807.03748}.

\bibitem[Radford et~al.(2021)]{radford2021learning}
Radford, A., Kim, J.~W., Hallacy, C., Ramesh, A., Goh, G., Agarwal, S.,
  Sastry, G., Askell, A., Mishkin, P., Clark, J., et~al. (2021).
\newblock Learning transferable visual models from natural language supervision.
\newblock In \textit{Proceedings of the International Conference on Machine
  Learning (ICML)}, pp.~8748--8763.

\bibitem[Robinson et~al.(2021)]{robinson2021contrastive}
Robinson, J., Chuang, C.-Y., Sra, S., and Jegelka, S. (2021).
\newblock Contrastive learning with hard negative samples.
\newblock In \textit{International Conference on Learning Representations
  (ICLR)}.

\bibitem[Zhang et~al.(2022)]{zhang2022contrastive}
Zhang, Y., Jiang, H., Miura, Y., Manning, C.~D., and Langlotz, C.~P. (2022).
\newblock Contrastive learning of medical visual representations from paired
  images and text.
\newblock In \textit{Machine Learning for Healthcare Conference (MLHC)},
  pp.~2--25.

\bibitem[Zhang et~al.(2023)]{zhang2023biomedclip}
Zhang, S., Xu, Y., Usuyama, N., Bagga, J.~K., Tinn, R., Preston, S.,
  Rao, R., Wei, M., Valluri, N., Wong, C., et~al. (2023).
\newblock Large-scale domain-specific pretraining for biomedical
  vision--language processing.
\newblock \textit{arXiv preprint arXiv:2303.00915}.

\end{thebibliography}

%% ---------------------------------------------------------
%% PHỤ LỤC
%% ---------------------------------------------------------
\appendix

\section{Cấu Hình Đầy Đủ Run V8}
\label{app:config}

\begin{table}[H]
  \centering
  \caption{Toàn bộ cấu hình siêu tham số của run v8 chính
    (\texttt{v8\_curriculum\_a100\_150ep\_seed42\_bs32\_v1}).}
  \label{tab:full_config}
  \small
  \begin{tabular}{ll}
    \toprule
    \textbf{Tham số} & \textbf{Giá trị} \\
    \midrule
    \multicolumn{2}{l}{\textit{Cơ bản}} \\
    epochs & 150 \\
    batch\_size & 32 \\
    grad\_accum & 4 \\
    effective\_batch & 128 \\
    freeze\_ep & 5 \\
    eval\_every & 5 \\
    seed & 42 \\
    num\_workers & 4 \\
    \midrule
    \multicolumn{2}{l}{\textit{Learning rate}} \\
    lr\_head (projection head) & 2e-4 \\
    lr\_enc (image encoder) & 1e-5 \\
    lr\_text (text encoder) & 1e-4 \\
    weight\_decay & 0.01 \\
    max\_grad\_norm & 1.0 \\
    \midrule
    \multicolumn{2}{l}{\textit{Clustering}} \\
    cluster\_start & 12 \\
    CLUSTER\_ALPHA & 0.35 \\
    proto\_start & 999 (vô hiệu hóa) \\
    rank\_start & 999 (vô hiệu hóa) \\
    mine\_start & 999 (vô hiệu hóa) \\
    \midrule
    \multicolumn{2}{l}{\textit{Clinical curriculum}} \\
    clinical\_start & 15 \\
    clinical\_ramp\_start\_weight & 0.03 \\
    clinical\_peak\_weight & 0.12 \\
    clinical\_ramp\_end & 25 \\
    clinical\_hold\_end & 60 \\
    clinical\_mid\_weight & 0.08 \\
    clinical\_decay\_end & 105 \\
    clinical\_final\_weight & 0.04 \\
    clinical\_best\_min\_strict & 3.5 \\
    \midrule
    \multicolumn{2}{l}{\textit{Trọng số loss}} \\
    aux\_weight & 0.02 \\
    rank\_weight & 0.0 \\
    \bottomrule
  \end{tabular}
\end{table}

\section{Phân Bố Trọng Số IDF Theo Nhóm Bệnh}
\label{app:idf}

Bảng~\ref{tab:idf_weights} tính trọng số IDF theo công thức~\eqref{eq:idf_weight}
từ tần suất trong Bảng~\ref{tab:label_dist} (bỏ ``No Finding'' và gộp
``Enlarged Cardiomediastinum'' vào ``Cardiomegaly'').

\begin{table}[H]
  \centering
  \caption{Trọng số IDF tương đối của 12 nhóm bệnh trong vector bệnh lý.
    Bệnh hiếm gặp (Pneumothorax, Fracture, Pleural Other) có trọng số cao hơn
    bệnh phổ biến (Cardiomegaly, Lung Lesion).}
  \label{tab:idf_weights}
  \begin{tabular}{lrrl}
    \toprule
    \textbf{Nhóm bệnh} & \textbf{Tần suất (\%)} & \textbf{Trọng số IDF (tương đối)} & \\
    \midrule
    Pleural Other           & 0.89 & Cao nhất  & $\bullet\bullet\bullet\bullet\bullet$ \\
    Fracture                & 1.82 & Rất cao   & $\bullet\bullet\bullet\bullet$ \\
    Pneumothorax            & 2.03 & Rất cao   & $\bullet\bullet\bullet\bullet$ \\
    Pneumonia               & 2.24 & Cao       & $\bullet\bullet\bullet$ \\
    Edema                   & 2.53 & Cao       & $\bullet\bullet\bullet$ \\
    Support Devices         & 3.06 & Vừa       & $\bullet\bullet$ \\
    Consolidation           & 5.49 & Vừa       & $\bullet\bullet$ \\
    Pleural Effusion        & 5.67 & Vừa       & $\bullet\bullet$ \\
    Cardiomegaly            & 9.00 & Thấp      & $\bullet$ \\
    Atelectasis             & 10.00 & Thấp     & $\bullet$ \\
    Lung Opacity            & 12.87 & Thấp     & $\bullet$ \\
    Lung Lesion             & 13.85 & Thấp nhất & $\cdot$ \\
    \bottomrule
  \end{tabular}
\end{table}

\section{Bộ Chỉ Số Đầy Đủ Tất Cả Checkpoint}
\label{app:all_checkpoints}

\begin{table}[H]
  \centering
  \caption{So sánh ba loại checkpoint của run v8 trên tập kiểm tra.}
  \label{tab:ckpt_compare}
  \begin{tabular}{l ccc c}
    \toprule
    \textbf{Checkpoint} & \textbf{Epoch} &
    \textbf{Strict Mean R@1} &
    \textbf{Cluster Mean R@1} &
    \textbf{Clinical-Valid Mean R@1} \\
    \midrule
    \texttt{best.pt}               & 120 & \textbf{3.85\%} & \textbf{69.36\%} & 59.86\% \\
    \texttt{best\_balanced.pt}     & 120 & 3.85\% & 69.36\% & 59.86\% \\
    \texttt{best\_clinical\_valid.pt} & 90 & 2.39\% & 65.78\% & 59.52\% \\
    \bottomrule
  \end{tabular}
  \vspace{0.3em}\\
  \small \textit{best.pt và best\_balanced.pt trùng kết quả do cùng epoch 120.
    best\_clinical\_valid.pt (epoch 90) cho strict thấp hơn rõ rệt trong khi
    clinical-valid chỉ tăng nhẹ (+0.34 pp), nên không được chọn làm checkpoint
    chính.}
\end{table}

\end{document}
